\usepackage{pdflscape}
\usepackage{fvextra}
\DefineVerbatimEnvironment{Highlighting}{Verbatim}{breaklines,commandchars=\\\{\},fontsize=\small}
\usepackage{scrlayer-scrpage}
\pagestyle{scrheadings}
\clearpairofpagestyles
\automark[chapter]{chapter}
\ihead{\rightmark}
\ofoot{\thepage}
\AtBeginDocument{
  \subtitle{Departament d'Arquitectura de Computadors \linebreak \linebreak Grau en Enginyeria Informàtica \linebreak Facultat d'informàtica de Barcelona \linebreak Universitat politècnica de Catalunya}
  \date{\today}
  \counterwithin{quartocalloutimpno}{chapter} % Comptadors de les callouts per capítol (RISC-V 4.1, EC 2.3, etc.) en lloc de global (RISC-V 1, RISC-V 2, ..., EC 1, EC 2, etc.)
  \counterwithin{quartocalloutnteno}{chapter}
  \counterwithin{quartocalloutwrnno}{chapter}
  \counterwithin{quartocallouttipno}{chapter}
  \counterwithin{quartocalloutcauno}{chapter}
  \definecolor{quarto-callout-note-color}{HTML}{198754} % mateixos colors HTML
  \definecolor{quarto-callout-note-color-frame}{HTML}{198754}
  \definecolor{quarto-callout-tip-color}{HTML}{ebca11}
  \definecolor{quarto-callout-tip-color-frame}{HTML}{ebca11}
  \definecolor{quarto-callout-caution-color}{HTML}{c0392b}
  \definecolor{quarto-callout-caution-color-frame}{HTML}{c0392b}
  \definecolor{quarto-callout-warning-color}{HTML}{8e44ad}
  \definecolor{quarto-callout-warning-color-frame}{HTML}{8e44ad}
  \definecolor{quarto-callout-important-color}{HTML}{0539f6}
  \definecolor{quarto-callout-important-color-frame}{HTML}{0539f6}
  \counterwithin{chapter}{part}
  \renewcommand{\thechapter}{\arabic{chapter}}  % Capítols numerats com 1, 2, 3, ... en lloc de 1.1, 1.2, 2.1, etc.
  \renewcommand{\chapterformat}{}               % No mostra el número del capítol
%  \renewcommand{\chapterpagestyle}{scrheadings} % E
  \renewcommand{\chapterpagestyle}{scrplain}
  \renewcommand{\chaptermarkformat}{}
}

% Algorismes (#thm-) en rodona (13_contrib.qmd §Blocs de codi). Quarto declara l'entorn
% theorem («Algorisme») amb \theoremstyle{plain}, que posa el cos en cursiva, i ho fa
% després d'aquest preàmbul. amsthm resol l'estil pel nom (\th@plain) a cada ús, de
% manera que n'hi ha prou de redefinir-lo en començar el document. Al llibre només
% theorem fa servir l'estil plain.
\makeatletter
\AtBeginDocument{\def\th@plain{\normalfont}}
\makeatother

% Etiqueta dels entorns #exr-: «Problema» a Problemes i Solucions, «Exercici» al laboratori
% (13_contrib.qmd §Problemari i solucionari). _quarto.yml dona «⁂» com a títol i prefix de
% l'entorn al PDF, i aquí «⁂» s'expandeix al nom vigent, que E1.qmd i L1.qmd canvien amb
% \renewcommand. No es pot posar \exercisename directament: Quarto n'escapa la contrabarra.
\usepackage{newunicodechar}
\newcommand{\exercisename}{Exercici}
\newunicodechar{⁂}{\exercisename}
